%% kapa: exam=1 round=27 year=2016 period=1 form=A subject=economics subject_ko=경제학원론 seq=34 number=74
%% source: Q-Net 감정평가사 기출문제 (cst00309) artlSeq=5169402 file=extracted/1008/5169402/제27회 감정평가사 1차 문제/제27회 감정평가사 1교시문제(A형).pdf
%% conversion: kordoc 4.24.51 -> tools/gen1.py (auto); review=proofread-sonnet; flags=none
\begin{question}{74}
실업에 관한 설명으로 옳은 것은?
\begin{choices}
  \item 만 15세 미만 인구도 실업률 측정 대상에 포함된다.
  \item 마찰적 실업은 자연실업률 측정에 포함되지 않는다.
  \item 더 좋은 직장을 구하기 위해 잠시 직장을 그만둔 경우는 경기적 실업에 해당한다.
  \item 경기적 실업은 자연실업률 측정에 포함된다.
  \item 현재의 실업률에서 실망실업자(discouraged workers)가 많아지면 실업률은 하락한다.
\end{choices}
\end{question}
\answer{5}
